\section{Introduction}

Weak random bonds in the planar Ising model provide a basic example in
which disorder is expected to change critical correlations through logarithmic
factors. The fluctuations of a quenched correlation are particularly revealing:
the correlation is first computed in a fixed environment, with its own Gibbs
normalization, and only then averaged over the bonds. This paper establishes
the predicted logarithmic growth of its relative second moment for small,
fixed binary disorder. The proof uses deterministic critical-reference
estimates from~\cite{ReferenceIsing}; the stochastic estimates and the
identification of the physical critical point are proved here.

\subsection{Model and result}

Let $E$ be the unordered nearest-neighbor edges of $\Z^2$, each counted once,
and let $(\xi_e)_{e\in E}$ be independent fair signs on a single probability
space. Fix $0<\eps<1$ and set $J_e(\omega)=1+\eps\xi_e(\omega)$.
For $\Lambda_n=\{-n,\ldots,n\}^2$, define
\begin{align}
 H^{\mathrm f}_{n,\omega}(\sigma)
   &=-\sum_{\substack{e=\{x,y\}\in E\\x,y\in\Lambda_n}}
       J_e(\omega)\sigma_x\sigma_y,\label{eq:free-hamiltonian}\\
 H^+_{n,\omega}(\sigma)
   &=-\sum_{\substack{e=\{x,y\}\in E\\e\cap\Lambda_n\ne\varnothing}}
       J_e(\omega)\sigma_x\sigma_y,
       \qquad \sigma_x=1\quad(x\notin\Lambda_n).
       \label{eq:plus-hamiltonian}
\end{align}
In the second sum, edges wholly outside $\Lambda_n$ are omitted. For
$b\in\{\mathrm f,+\}$ and $\beta\ge0$, write
$\langle\cdot\rangle^b_{n,\beta,\omega}$ for the Gibbs expectation
proportional to $\exp(-\beta H^b_{n,\omega})$. Every expectation is
normalized separately at its indicated environment and boundary condition.

The physical critical inverse temperature is
\begin{equation}\label{eq:critical-definition}
 m_\eps(\beta)=\E\!\left[\lim_{n\to\infty}
       \langle\sigma_0\rangle^+_{n,\beta,\omega}\right],
 \qquad
 \beta_c(\eps)=\inf\{\beta\ge0:m_\eps(\beta)>0\}.
\end{equation}
Ferromagnetic monotonicity gives the plus limit. Comparison with constant
bonds, with $\beta_0=\tfrac12\log(1+\sqrt2)$, gives
\[
 0<\frac{\beta_0}{1+\eps}\le\beta_c(\eps)
       \le\frac{\beta_0}{1-\eps}<\infty.
\]
For each integer $r\ge3$, let
\begin{equation}\label{eq:quenched-correlation}
 C_{\eps,\omega}(r)=\lim_{n\to\infty,\,n\ge r}
       \langle\sigma_0\sigma_{(r,0)}\rangle^{\mathrm f}
                   _{n,\beta_c(\eps),\omega},
 \qquad
 R_\eps(r)=\frac{\E[C_{\eps,\omega}(r)^2]}
                 {\E[C_{\eps,\omega}(r)]^2}.
\end{equation}
The free correlation limit exists by monotonicity in the volume and lies
in $(0,1]$. In particular $R_\eps(r)$ is defined and is at least one.
All unqualified expectations $\E$ below are over the environment; auxiliary
independent random choices used in the proof will be included explicitly
when introduced.

\begin{theorem}\label{thm:main}
Assume the deterministic reference-model estimates of
\cite{ReferenceIsing}, at the scopes stated in Section~\ref{sec:reference}.
There is $\eps_0>0$ such that for every fixed $\eps\in(0,\eps_0)$,
\begin{equation}\label{eq:main-limit}
 \lim_{r\to\infty,\,r\in\N}
       \frac{\log R_\eps(r)}{\log\log r}=\frac14.
\end{equation}
The critical inverse temperature in \eqref{eq:quenched-correlation} is
the physical value \eqref{eq:critical-definition}; it is the unique positive
solution of
\begin{equation}\label{eq:critical-curve}
 \sinh\bigl(2\beta(1-\eps)\bigr)
 \sinh\bigl(2\beta(1+\eps)\bigr)=1.
\end{equation}
\end{theorem}

The order of limits is essential. The disorder strength is fixed first;
the free thermodynamic limit is then taken for each separation, followed
by the disorder moments and the limit in $r$. Constants and errors may
depend on this fixed positive strength. At $\eps=0$ the environment is
deterministic and $R_0(r)=1$.

The reference input concerns deterministic translation-invariant,
square-symmetric, finite-range even perturbations and their critical
nearest-neighbor comparison kernels. It supplies neither a random critical
point nor a quenched moment estimate. Section~\ref{sec:reference} states
the particular reference consequences needed below, including their
buffer and symmetry conditions. The proof of
Theorem~\ref{thm:main} establishes the random estimates from these inputs.

\subsection{Historical context}

The pure planar Ising model has specific-heat exponent zero, the borderline
case left undecided by Harris's criterion for the effect of quenched
disorder~\cite[Section~5]{Harris1974}. Dotsenko and Dotsenko obtained the
double-logarithmic prediction for the random model's specific heat
\cite{DotsenkoDotsenko1983}. Shalaev's renormalization analysis predicted
unchanged leading critical indices~\cite{Shalaev1984}. Shankar and Ludwig
developed the logarithmic corrections to spin-correlation moments
\cite{Shankar1987,Ludwig1990}. In particular, Shankar predicted
the factor $(\log r)^{1/4}$ in the second moment of the spin correlation,
and Ludwig studied logarithmic corrections to its moments. The exponent
in Theorem~\ref{thm:main} is thus an established prediction. The issue
addressed here is its derivation for normalized quenched correlations at
the physical critical point, under the stated reference input.

The deterministic reference theory has a different role. Critical
nearest-neighbor spin correlations and mixed primary fields have rigorous
scaling limits; the limiting fields satisfy operator-product expansions
\cite{CHI2015,CHI2023}. Rigorous results for deterministic weak
perturbations include energy correlations in the plane
\cite{GGM2012}, energy correlations on cylinders~\cite{AGG2023}, and
boundary spin correlations in the half-plane~\cite{CGG2025}. These results
provide context for the deterministic bulk premise used here
\cite{ReferenceIsing}; none supplies the random moment estimates below.
The energy and spin operator products enter the coefficient computation
in Section~\ref{sec:coefficients}, after the local estimates needed for
random interactions have been established.

Rigorous critical behavior at a self-dual random-bond law also predates
this work. Building on disorder-rounding results of Aizenman and Wehr
\cite{AizenmanWehr1990}, Chayes and Shtengel proved zero magnetization
and algebraic lower bounds on averaged correlations at the binary
self-dual law~\cite[Theorems~1$'$ and~2$'$]{ChayesShtengel1999}. Their
argument did not exclude a critical phase containing that point. Here
the weak-disorder estimates identify the tuned parameter with the
physical critical threshold and give the logarithmic moment asymptotic
there. The probabilistic part uses the random-cluster representation
\cite{FK1972,EdwardsSokal1988}, critical crossing estimates
\cite{DHN2011,CDH2016}, product-space influences~\cite{GJN2013}, and the
monotone-measure extension~\cite{DCRT2019} of the decision-tree inequality
of O'Donnell, Saks, Schramm, and Servedio~\cite{OSSS2005}.

\subsection{Proof strategy and the role of locality}

It is convenient to include inverse temperature in the couplings and use
the coordinates
\begin{equation}\label{eq:model-td}
 \log\sinh(2K_e)=t+D\xi_e,\qquad D>0.
\end{equation}
For each sufficiently small fixed $D$, we first choose $t$ by a
renormalization construction. The last part of the proof identifies this
choice with $t=0$ and then returns to the original parameter $\beta$.

At scale $r_j=r_0L^j$, the effective interaction is a sum of local even
spin functions. Each cell has one distinguished term whose size is allowed
to be large. The other terms are indexed by connected sets of cells and
have small norms with exponential decay in the size of those sets. Every
term records both the spins on which it acts and the original random
variables on which its coefficients depend. This second record will give
exact independence of the leading factors attached to well-separated
spin insertions.

The first difficulty is to transport interactions that are small in high
moments but are not uniformly small. Section~\ref{sec:blocking} constructs
an exact local change of representation. An isolated group of large
interactions is integrated as a positive exponential weight. Insertions
in the same region are transported by its normalized conditional law.
Reference comparison across a surrounding buffer bounds the resulting
density, and reference even-function decay attenuates its dependence on
the outer spins. Failures of this attenuation have a distant large-field
witness, which can be paid for using independence and high moments.
Section~\ref{sec:norms} turns these observations into the estimates needed
to iterate the construction. In particular, it controls mixed moments of
spin sources and nearby background fields, whose randomness is shared.

The linearized background transport has one expanding mean coordinate
and one neutral covariance coordinate. The first is removed by choosing
$t$; the second, denoted $g_j$, measures thermal-charge variance per unit
area. Section~\ref{sec:linear} proves contraction of the remaining
coordinates. The nonlinear coefficients of the two surviving effects are
computed in Section~\ref{sec:coefficients}. Its method uses finitely many
microscopic energy insertions: exact partition identities preserve their
normalized charges, and the reference energy and spin operator products
determine the logarithmic terms. Only finite-degree polynomial calculations
are needed. This separates the coefficient calculation from the estimates
that justify the real random iteration.

The resulting flow, proved in Sections~\ref{sec:flow} and~\ref{sec:marked},
satisfies $g_j\sim(B_0j)^{-1}$ for a positive constant $B_0$. The leading
random factor of a transported spin has a bounded positive mean and
second moment $j^{1/8+o(1)}$. Two distant insertions have independent
leading factors once their small nonlocal remainders are removed.
Section~\ref{sec:readout} uses an enlarging terminal buffer to compare
the remaining Gibbs expectation with the critical reference expectation.
It obtains the relative second-moment exponent $1/4$ and also a high-moment
wired-arm bound.

Finally, Section~\ref{sec:transition} identifies the physical transition.
Twisted torus partition functions locate the tuned value at the
distributionally self-dual point. A product-space influence inequality
and a random-shell exploration inequality make the needed threshold and
subcritical estimates explicit. Planar random-cluster duality then gives
positive magnetization above tuning. This final step connects the local
renormalization construction to the critical point used in the statement.

The reusable parts of the method are the positive treatment of large
local interactions, the simultaneous tracking of spin support and random
dependence, and the extraction of marginal coefficients from invariant
charges of finitely many insertions. Their application here requires all
three: locality controls dependence, high moments control exceptional
fields, and the coefficient identities determine the logarithmic exponent.
